% =====================================================================
\section{Project Set-up and How to Reproduce Everything}
\label{sec:setup}
% =====================================================================

\subsection{Repository map}

\begin{center}\small
\begin{tabularx}{\linewidth}{@{}>{\raggedright\arraybackslash}p{5.4cm} X@{}}
\toprule
\textbf{Path} & \textbf{What it is}\\
\midrule
\code{base paper.pdf} & the base paper\\
\code{paper\_reproduction.py} & reproduction of the paper's Section IV (Part 1, Part 2), sensitivity, Monte-Carlo; writes \code{docs/figures/fig*.png}, \code{docs/results/metrics.*}\\
\code{x3\_sim/x3\_model.py} & software-in-the-loop surrogate of the Gazebo X3 (all constants from the Gazebo model files)\\
\code{x3\_sim/controllers.py} & \code{OriginalGazeboController} (exact port of our first node) and \code{FixedController} (the corrected paper controller)\\
\code{x3\_sim/run\_x3\_experiments.py} & experiments E0--E5 on the X3; writes \code{docs/figures/x3\_*.png}, \code{docs/results/x3\_metrics.*}\\
\code{gazebo.py} & ROS~2 node (v2) that flies the Gazebo X3 with \code{FixedController}\\
\code{plot\_gazebo\_log.py} & turns a log of \code{gazebo.py} into figures and metrics\\
\code{tests/test\_gazebo\_node\_offline.py} & runs the real \code{gazebo.py} without ROS against the surrogate\\
\code{toy\_model/} & MATLAB toy model, steps 1--8, \code{run\_all.m}, \code{lib/}\\
\code{Drones\_Mujoco\_*.ipynb}, \code{One\_drone\_alligning.ipynb} & MuJoCo 3-D attitude and flight demos\\
\code{index.html}, \code{app.js}, \code{simulation.js}, \code{style.css} & interactive browser simulator (illustration)\\
\code{legacy/} & first versions kept for the record (\code{gazebo\_v1.py}, early drafts); not used as evidence\\
\code{docs/} & this document, the mathematical models, all figures and results\\
\bottomrule
\end{tabularx}
\end{center}

\subsection{Tools}

\begin{center}\small
\begin{tabularx}{\linewidth}{@{}l X@{}}
\toprule
\textbf{Tool} & \textbf{Used for}\\
\midrule
Python 3 + NumPy + Matplotlib & reproduction, X3 surrogate, plots (no other dependency)\\
MATLAB (no toolbox) & toy model and independent cross-check of the reproduction\\
ROS~2 + Gazebo Sim (gz) & X3 quadcopter world; \code{ros\_gz\_bridge} for \code{/world/quadcopter/dynamic\_pose/info} and \code{/X3/gazebo/command/motor\_speed}\\
MuJoCo (Python) & 3-D visual demos of the attitude controller\\
\LaTeX\ (\code{pdflatex}) & both documents\\
\bottomrule
\end{tabularx}
\end{center}

\subsection{Commands}

\begin{lstlisting}[style=sh]
pip install numpy matplotlib
python3 paper_reproduction.py            # Part 1, Part 2, sensitivity, Monte-Carlo  (~6 min)
python3 x3_sim/run_x3_experiments.py     # X3 experiments E0-E5                      (~16 min on 2 cores)
python3 tests/test_gazebo_node_offline.py  # runs gazebo.py against the X3 surrogate (~2 min)
\end{lstlisting}
In MATLAB, from any folder:
\begin{lstlisting}[style=mat]
run('<repo>/toy_model/run_all.m')
\end{lstlisting}
It writes \code{toy\_model/figures/*.png} and the
complete printed output to \code{toy\_model/results\_log.txt}.

Gazebo (with the \code{quadcopter} world running and the two topics bridged):
\begin{lstlisting}[style=sh]
python3 gazebo.py --ros-args -p controller_mode:=constrained -p experiment:=hover -p duration:=180.0
python3 gazebo.py --ros-args -p experiment:=square -p duration:=45.0
python3 gazebo.py --ros-args -p controller_mode:=pd -p disturbance_scale:=0.5
python3 plot_gazebo_log.py /tmp/drone_constrained_hover.npz
\end{lstlisting}

\begin{projbox}[reproducibility]
Running \code{paper\_reproduction.py} again on a clean checkout reproduced \code{docs/results/metrics.md} and
\code{metrics.json} \emph{bit for bit} (the Monte-Carlo uses a fixed seed, 7). The X3 experiments are deterministic as well
(fixed seeds), and the offline test of \code{gazebo.py} passes (hover ends at $z=2.000$\,m, the square returns to
$(0,0,2)$\,m).
\end{projbox}

\subsection{Simulation settings common to all results}

\begin{center}\small
\begin{tabularx}{\linewidth}{@{}l X@{}}
\toprule
\textbf{Setting} & \textbf{Value}\\
\midrule
Integrator & fixed-step RK4, $\Delta t=1$\,ms; quaternion re-normalised after every step\\
Control & computed once per step and held (zero-order hold), as in a digital controller\\
Saturation & Eq.~(3)--(4), applied to the commanded torque\\
Barrier & Eq.~(22) (boundary layer), $\varepsilon$ fed to the barrier limited to 0.999 (never active inside the envelope)\\
Adaptive law & explicit Euler, $\that\ge0$\\
Unstated parameters & $c=0.5$, $\varsigma=0.02$, Part~2: $T_s=15$\,s, $\rho_{q0}=0.36$ (Mathematical Models, Section 2)\\
Horizon & 20\,s (the paper's plots)\\
\bottomrule
\end{tabularx}
\end{center}

\subsection{Metrics (defined once, used everywhere)}
\label{sec:metrics}

With $T$ the simulation length and $\mathbb 1[\cdot]$ the indicator function:
\[
V_q=\frac{100}{T}\int_0^T\mathbb 1\big[\exists i:|q_i|\ge\rho_q\big]dt,\qquad
V_\omega=\frac{100}{T}\int_0^T\mathbb 1\big[\exists i:|\omega_i|\ge\rho_\omega\big]dt\qquad(\text{\% of time outside}),
\]
\[
\varepsilon_{\max}=\max_{t,i}|\varepsilon_{li}(t)|,\qquad
S_u=\frac{100}{T}\int_0^T\mathbb 1\big[\exists i:|u_i|\ge u_{\max}\big]dt,\qquad
E_u=\int_0^T\|\sat(\vu)\|^2dt,
\]
plus the settling time (last instant with $\|\qv\|>0.01$), the final $\|\qv\|$ and $\|\vw\|$. A run is
\emph{violation-free} if $V_q=V_\omega=0$.
